\SourceLeaf{14}
Dabei gelten die oberen oder die unteren Vorzeichen, je nachdem, ob $M$ außerhalb oder innerhalb des Körpers liegt. Im ersten Fall bilden die Kegelabschnitte von $P'$ bis $P''$, von $P'''$ bis $P''''$ und so weiter einen Teil des festen Körpers, im zweiten die Abschnitte von $M$ bis $P'$, von $P''$ bis $P'''$ und so weiter. In beiden Fällen ergibt sich daher für den Teil des Körpers, der im Inneren des Kegels auf der Grundfläche $\mathrm{d}\Sigma$ liegt
\[
=-\tfrac13(r'\mathrm{d}s'.\cos MQ'+r''\mathrm{d}s''.\cos MQ''+r'''\mathrm{d}s'''.\cos MQ'''+\text{ etc.})
\]
Behandeln wir alle Elemente $\mathrm{d}\Sigma$ auf dieselbe Weise und summieren, so erhalten wir
\par\medskip\Needspace{4\baselineskip}
\begin{center}FÜNFTER LEHRSATZ.\end{center}
\par
\textit{Das Volumen des ganzen Körpers ist gleich dem über die gesamte Oberfläche des Körpers erstreckten Integral $-\tfrac13\int r\mathrm{d}s.\cos MQ$.}

\Article{8}
\par
Nehmen wir nun an, der Körper habe gleichmäßige Dichte und jedes seiner Elemente übe auf den Punkt $M$ eine Anziehung aus, die einer Funktion des Abstands proportional ist. Bezeichnet $\rho$ den Abstand des Elements vom angezogenen Punkt, so wird die Anziehung durch das mit $f\rho$ multiplizierte Volumen des Elements ausgedrückt. Man denke sich zunächst unseren Kegel auf der Grundfläche $\mathrm{d}\Sigma$ vollständig mit Materie gefüllt und durch Kugeloberflächen, die einander unendlich nahe sind und um den Mittelpunkt $M$ beschrieben werden, in unendlich viele Elemente zerlegt. Ein solches Element an der Kugel mit Radius $=\rho$ wird durch $\rho\rho\mathrm{d}\rho.\mathrm{d}\Sigma$ ausgedrückt, also die Kraft, mit der es auf $M$ wirkt, durch $\mathrm{d}\Sigma.\rho\rho f\rho.\mathrm{d}\rho$. Bezeichnen wir daher das Integral $\int\rho\rho f\rho.\mathrm{d}\rho$ mit $\Phi\rho$, so ist klar, dass $\mathrm{d}\Sigma.(\Phi\rho-\Phi0)$ die Anziehung des vom Scheitel bis zum Abstand $\rho$ reichenden Kegelabschnitts auf den Punkt $M$ ausdrückt oder, allgemeiner, dass $\mathrm{d}\Sigma.(\Phi\rho'-\Phi\rho)$ die Anziehung des Kegels zwischen den Abständen $\rho$ und $\rho'$ vom Scheitel ausdrückt. Alle innerhalb des Kegels liegenden Teile unseres Körpers ziehen daher den Punkt $M$ in Richtung $M\Pi$ mit einer Kraft an, die ausgedrückt wird durch
\[
\mathrm{d}\Sigma.(-\Phi r'+\Phi r''-\Phi r'''+\text{ etc.})
\]
wenn $M$ außerhalb des Körpers liegt, oder durch
\[
\mathrm{d}\Sigma.(-\Phi0+\Phi r'-\Phi r''+\Phi r'''-\text{ etc.})
\]
wenn $M$ innerhalb des Körpers liegt, oder gleichwertig
\par
im ersten Fall durch
